\section{The Climb, From Field to Mind}

The physics is the bottom of a ladder, not the whole story. The claim of this part is that the same four base things that build the field also build life, mind, and intelligence, as higher regimes of one substrate rather than as new ingredients. This part climbs that ladder. It is more speculative than the chapters on physics, and it works at toy scale rather than at the scale of a real organism or brain. We mark each rung's status plainly, and the reader should treat the upper rungs as plausibility arguments backed by separate demonstrations rather than as a single end-to-end proof.

\subsection{Everything is a glider}

In the rule, a thing is not a fixed set of cells. It is a pattern that persists while the cells beneath it come and go, the way a \emph{glider} in a cellular automaton re-forms one step over with each beat. Every particle and every self is this kind of object, a stable traveling pattern made of nothing but arrangement. This is why a self can persist while none of its underlying cells do. In the carrying image, a flame keeps its shape while the gas burning in it is always new. Holding this picture in mind is the key to the whole climb, because it explains how identity can survive complete material turnover at every level, from a particle to a person.

\subsection{The multiscale shortcut}

We do not simulate a mind from the base. A vibe is sub-Planck, a particle is on the order of $10^{60}$ vibes, and a brain is hopeless to build cell by cell. Instead each layer has an up-map that coarse-grains a block into one richer symbol, and a down-map that lifts a symbol to a representative pattern. A layer is faithful when its dynamics commute, so that coarse-graining then evolving equals evolving then coarse-graining. This was checked across scales from thirty thousand to one million cells \cite{pollard2026vibetest}. So we study each rung at its own level, the way biology studies cells without simulating quarks.

The maps have a precise shape, a Galois-like connection rather than an isomorphism. The up-map is lossy and lands in a richer macro-alphabet, while the down-map is a canonical lift, with the composite returning the symbol exactly and the dynamics square commuting. The key consequence is that the base never limits the macro-alphabet, only the macro-dynamics. Composites can be arbitrarily rich, the way bits become files, so the objection that the tone is only ternary has no force at high layers. The ternary base constrains how things move, not how much they can encode. A higher vibe is an aggregate read-only view of the one stored base. There is no hidden super-tone, which is the no-hidden-state commitment in which the note is perception.

\begin{principle}[Nothing added up the ladder]
Each higher layer is a faithful coarse-graining of the layer below it, with a lossy up-map into a richer alphabet and a canonical down-map back. The base limits only the dynamics, never the alphabet, so composites can be arbitrarily rich while remaining a read-only view of the one stored base. No new ingredient is introduced at any rung.
\end{principle}

\subsection{What the middle layers are}

Raw tones are not particles. There are one to a few emergent middle layers between the vibe churn and physics, and we can name the first one. The first faithful middle layer is a slow conserved-charge density mode, a hydrodynamic slow-mode whose autocorrelation lifetime grows with coarse-graining scale. A region's charge changes only by boundary flux, so larger regions are slower. This is measured, with the lifetime climbing from near zero at the base to a finite value a few levels up, and it is explicitly not physics. It is the first slow collective variable above the churn, the way a fluid density emerges above molecules.

The caveats matter here. The lifetime is modest, and hyperbolic geometry weakens the slow mode, since a ball there is nearly all boundary. That is one more reason physics lives on the flat cusp rather than in the bulk. So the bottom of the ladder reads as tone churn, then a slow density mode, then on the flat cusp the relativistic field, then particles. At no point is a tone the same thing as a particle. Keeping this clear is important, because the temptation to read raw clustering as gravity or matter is exactly the kind of overreach the theory is built to avoid.

\subsection{The climb, with status}

\begin{itemize}
  \item \textbf{Field to particle, measured.} A particle is a persistent bump in the field. Its centroid moves in a straight line at constant speed below the light cone, an Ehrenfest result on the lattice with $R^2 \approx 1$.
  \item \textbf{Particle to composite, measured.} Two particles keep a uniform center of mass, with momentum conserved through interaction at $R^2 \approx 0.997$. The conserved charge gives the minimal atom, a net-charged region that cannot vanish.
  \item \textbf{Atoms through molecules, cells, tissues, senses, to minds, plausible.} This is not simulated as one run. Every ingredient is separately demonstrated, including binding by the merge law, self-maintenance, copying, the onion of senses, and the integrated hub. It is a grounded extrapolation, not a single proof.
\end{itemize}

The one-line version of the ladder runs as follows. A particle is a bump in the field. An atom is a conserved bound unit. A molecule is a larger bound unit. A cell is a self-maintaining self. A tissue is a tower of cells. A sense is a tuned boundary. A mind is an integrated hub with a self-model. A human is the high end of that same climb. None of these requires a new base thing, which is the entire point of laying them out as one continuous ascent.

\begin{result}[The lower rungs are measured]
The field-to-particle and particle-to-composite steps are measured directly, with straight-line centroid motion below the light cone at $R^2 \approx 1$ and momentum conservation through interaction at $R^2 \approx 0.997$ \cite{pollard2026vibetest}. The rungs above the composite are demonstrated ingredient by ingredient rather than in one run, and are marked plausible.
\end{result}

The substrate does not stop at physics. The same rule, read at greater heights, gives the living and the mental. This is the boldest and most speculative claim of the theory, and it is the reason this part is laid out as a ladder with the rungs plainly tagged. The chapters that follow take up life and then mind in turn, each one keeping the same discipline of separating what is measured from what is extrapolated.
